\subsection*{Maquetación del front-end}
\addcontentsline{toc}{subsection}{Maquetación del front-end}

\us[
	identificador=US-MF-01,
	titulo={Maquetación de la pantalla de inicio},
	como={Usuario},
	quiero={Entender la funcionalidad de la aplicación desde el momento del acceso},
	requerimientos={\reqref{RF-GU-01}, \reqref{RF-GU-02}, \reqref{RF-IR-07}, \reqref{RNF-CO-01}, \reqref{RNF-CO-02}, \reqref{RNF-US-01}, \reqref{RNF-US-03}, \reqref{RNF-RE-01}},
	puntos={8},
	criterios={
			\begin{itemize}
				\item La pantalla de inicio tiene una barra de navegación.
				\item La barra de navegación presenta el logo y un acceso directo a la pantalla de
				      inicio.
				\item La barra de inicio presenta opciones para registrarse o iniciar sesión.
				\item La pantalla de inicio tiene una hero section.
				\item La hero section describe claramente el propósito de la aplicación.
				\item La hero section presenta un \textit{call to action} que guía el viaje de
				      usuario hacia el registro o el inicio de sesión.
				\item La pantalla de inicio presenta una lista de recetas destacadas.
				\item La pantalla de inicio presenta un pie de página con información de contacto y
				      enlaces a los documentos legales.
			\end{itemize}
		},
	tareas={
			\begin{itemize}
				\item Selección de la paleta de colores, tipografías y elementos gráficos que
				      representarán la marca \textbf{Muncher} en todas las pantallas.
				\item Selección de tecnologías front-end a utilizar: librería de componentes
				      (preexistente o propia), sistema de gestión de estilos (CSS puro, CSS Modules,
				      CSS in JS o Tailwind CSS).
				\item Maquetación de la pantalla de inicio de forma responsiva y accesible.
				\item Implementación de la pantalla de inicio en el front-end.
			\end{itemize}
		}
]

\us[
	identificador=US-MF-02,
	titulo={Internacionalización},
	como={Usuario},
	quiero={Acceder al contenido en mi idioma preferido},
	requerimientos={\reqref{RF-GU-01}, \reqref{RF-GU-02}, \reqref{RF-IR-07}, \reqref{RNF-CO-01}, \reqref{RNF-CO-02}},
	puntos={8},
	criterios={
			\begin{itemize}
				\item Los textos de la aplicación se almacenan de forma separada al código.
				\item Los textos de la aplicación se encuentran en varios idiomas.
				\item El usuario puede seleccionar el idioma de la aplicación.
				\item La aplicación detecta el idioma del navegador del usuario y lo utiliza como
				      idioma por defecto.
			\end{itemize}
		},
	tareas={
			\begin{itemize}
				\item Implementación de la internacionalización en el front-end.
			\end{itemize}
		}
]
